% Revised draft (Phase 2), based on draft_final_review4.tex.
% Former title: "Evaluating Detection Rule Efficacy: A Reproducible Benchmark
% for Security Operations". The title was changed because no detection rule
% was executed in the reported runs. AUTHOR: approve or revert the title.
% Lines marked "AUTHOR:" need the author's decision before submission.
\documentclass[11pt,a4paper]{article}

\usepackage[utf8]{inputenc}
\usepackage[T1]{fontenc}
\usepackage{graphicx}
\usepackage{booktabs}
\usepackage{hyperref}
\usepackage{listings}
\usepackage{xcolor}
\usepackage{geometry}
\usepackage{amsmath}
\usepackage{float}

\geometry{margin=1in}

\definecolor{codegreen}{rgb}{0,0.6,0}
\definecolor{codegray}{rgb}{0.5,0.5,0.5}
\definecolor{codepurple}{rgb}{0.58,0,0.82}
\definecolor{backcolour}{rgb}{0.95,0.95,0.92}

\lstdefinestyle{mystyle}{
    backgroundcolor=\color{backcolour},
    commentstyle=\color{codegreen},
    keywordstyle=\color{magenta},
    numberstyle=\tiny\color{codegray},
    stringstyle=\color{codepurple},
    basicstyle=\ttfamily\footnotesize,
    breakatwhitespace=false,
    breaklines=true,
    captionpos=b,
    keepspaces=true,
    showspaces=false,
    showstringspaces=false,
    showtabs=false,
    tabsize=2
}
\lstset{style=mystyle}

\title{Telemetry Availability and Log-Availability Latency in a Containerized
Detection Benchmark: An Experience Report}
\author{SOC Detection Engineering Benchmark} % AUTHOR: replace with author name(s) and affiliation
\date{2026-10-10}

\begin{document}

\maketitle

\begin{abstract}
Benchmarks for security detection content depend on the test environment
producing the log evidence that a detection would need. We report what happened
when we built a small containerized Linux lab to exercise 11 MITRE ATT\&CK
techniques and measured, for each, whether log evidence reached Elasticsearch
within a collection window after the attack command ran. Evidence was matched
with keyword-based Elasticsearch queries; no Sigma rule or other detection rule
was executed. Evidence was observed for 2 of the 11 planned techniques, each in
a single run: local account creation (T1136.001), with a log-availability
latency of 3.20\,s, and cron persistence (T1053.003), with a latency of
30.94\,s, which required a 65\,s post-attack window after a 15\,s window
missed the event. The other 9 techniques were not evaluated after the harness
was corrected: the audit daemon, their expected log source, did not run in the
container, which we attribute to the Linux kernel permitting audit control only
from the initial PID namespace. No false-positive baseline was measured
and runs were not repeated, so the results are observations about this lab, not
estimates of detection performance. We describe the measurement problems we
encountered and what they imply for building such a benchmark.
\end{abstract}

\section{Introduction}

Detection engineers test detection content by running attack behaviors in a
lab and checking whether the expected alerts appear. That check has a
precondition that is easy to overlook: the lab must generate, ship, and index
the log events the detection reads, and must do so within the time the test
harness waits. If it does not, a missed detection says nothing about the
detection logic.

We built a containerized lab and harness intended to benchmark detection rules
against ATT\&CK-mapped attack commands. The work did not reach that goal. Only
2 of 11 planned techniques produced observable log evidence, and the evidence
was matched by keyword queries, not by rules. This paper reports that outcome
and the reasons for it.

\subsection{Research Question}

\begin{quote}
In a containerized Linux lab with file-based log shipping, which log sources
yielded observable evidence for a planned set of ATT\&CK techniques, and how
did source-specific event-generation delay affect the collection window needed
to observe that evidence?
\end{quote}

\subsection{Contributions}

\begin{itemize}
    \item An account of why 9 of 11 planned techniques could not be evaluated
    in this lab, attributed to a kernel restriction on audit control outside
    the initial PID namespace (Section~\ref{sec:telemetry}).
    \item A worked example of a collection-window failure: a 15\,s window
    missed cron evidence that a 65\,s window captured, because cron checks
    for crontab changes once per minute (Section~\ref{sec:cron}).
    \item A list of measurement-validity problems found while auditing our own
    results, including attribution confounds and a log-compression artifact
    (Section~\ref{sec:threats}).
    \item The harness, lab configuration, and evidence files, with a unit test
    suite of 60 tests covering the harness logic.
\end{itemize}

\subsection{Scope}

The paper does not evaluate detection rules. The repository contains 10 Sigma
rule files, but the harness never loaded them; all reported matches came from
keyword queries (Section~\ref{sec:query}). The paper reports no false-positive
rate and no repeated runs. We use \emph{observed} to mean that at least one
indexed log document matched a technique's keywords inside its collection
window.

\section{Background and Related Work}

\paragraph{ATT\&CK and adversary emulation.}
MITRE ATT\&CK is a knowledge base of adversary tactics and techniques
\cite{attack-design}. Automated emulation of ATT\&CK techniques dates at least
to Applebaum et al.\ \cite{applebaum2016}, and Atomic Red Team provides a
library of tests mapped to ATT\&CK \cite{atomic}. MITRE's own evaluations apply
emulation to commercial endpoint products \cite{attack-evals}; Shen et al.\
analyzed the results of those evaluations \cite{shen2024}. Virkud et al.\
examined ATT\&CK-annotated rule sets from Carbon Black, Splunk, Elastic, and
Sigma and concluded that technique coverage is a limited measure of security
posture, in part because many techniques cannot be realized as detection rules
\cite{virkud2024}. Our count of observed techniques should be read with that
caution in mind.

\paragraph{Rule-based detection and Sigma.}
Sigma is a generic signature format for SIEM systems \cite{sigma-spec}; rules
are converted to backend queries by tools such as pySigma \cite{pysigma}. Uetz
et al.\ showed that a substantial share of the Sigma rules they analyzed could
be evaded by small changes to attack commands \cite{uetz2024}. Recent preprints
study how public rule repositories change over time \cite{long2026} and how
rule sets can be optimized \cite{shukla2025}. These studies analyze or execute
real rules. Our keyword queries are much cruder, and our results are not
comparable with theirs.

\paragraph{Testbeds and log datasets.}
The closest prior work builds testbeds that generate log data reproducibly.
SOCBED is an open testbed designed so that log data generation can be repeated
and adapted \cite{uetz2021}. Landauer et al.\ generate labelled log datasets
from model-driven testbeds \cite{landauer2021,landauer2022}. These projects
operate at a larger scale than ours and repeat their scenarios; we ran each
technique once.

\paragraph{Alert quality.}
Axelsson showed that the false-alarm rate limits the usefulness of an
intrusion detector when attacks are rare \cite{axelsson2000}, which is why a
detection count without a false-positive measurement is incomplete. Alahmadi et
al.\ found that SOC practitioners attribute most false positives to benign
triggers, that is, alarms that fire correctly on legitimate activity
\cite{alahmadi2022}. Studies of SOC practice document the cost of rule
maintenance and alert handling \cite{vermeer2022,vermeer2023,kokulu2019}.

\paragraph{Host telemetry in containers.}
The Linux audit subsystem is not namespace-aware, which complicates auditing
of containers \cite{lim2021,chen2021}. Houssel et al.\ compare kernel telemetry
options and identify eBPF as the most promising \cite{houssel2026}. The
specific restriction that affected our lab is in the kernel source
\cite{linux-audit} and is described in Section~\ref{sec:telemetry}.

\paragraph{Evaluation methodology.}
McHugh's critique of the DARPA intrusion detection evaluations
\cite{mchugh2000}, the survey of benchmarking flaws by van der Kouwe et al.\
\cite{vanderkouwe2019}, and Collberg and Proebsting's study of repeatability
\cite{collberg2016} motivate the explicit reporting of scope, sample size, and
failed runs in this paper.

\paragraph{Gap.}
We did not find prior work that reports, for a small containerized lab, which
log sources were available per technique and how event-generation timing
constrained the collection window. Our search was a narrative review, not a
systematic one, so we do not claim that no such work exists.

\section{Lab and Measurement Method}

\subsection{Lab Architecture}

The lab was a Docker Compose stack on a macOS host:

\begin{enumerate}
    \item \textbf{Target container}: Ubuntu 22.04, run with
    \texttt{privileged: true}, with rsyslog, cron, and the auditd package
    installed.
    \item \textbf{Log shipper}: Filebeat 8.12.0, reading the target's
    \texttt{/var/log} through a shared volume.
    \item \textbf{Log store}: Elasticsearch 8.12.0 (Kibana 8.12.0 was also
    deployed but not used in the measurements).
    \item \textbf{Harness}: Python code on the host that ran attack commands
    with \texttt{docker exec} and then queried Elasticsearch.
\end{enumerate}

\subsection{Technique Set}

The attack manifest listed 11 ATT\&CK techniques with 12 test commands (T1105
had two). The commands were written by hand for this project. The target image
clones the Atomic Red Team repository, but the harness did not execute Atomic
Red Team tests.

\subsection{Telemetry Availability}
\label{sec:telemetry}

The manifest named auditd as the expected log source for all 11 techniques.
The audit daemon did not start in the target container. An earlier draft of
this report recorded the error \texttt{Error setting audit daemon pid
(Operation not permitted)}; the error output was not saved in the results
directory. A read-only status query in the running container on 2026-10-13
found no audit daemon process, and \texttt{auditctl -s} returned
\texttt{Error sending status request (Operation not permitted)}. We attribute
the failure to a restriction in the kernel source: audit control
messages are rejected with \texttt{EPERM} unless the sender is in the initial
PID namespace \cite{linux-audit}. The container had its own PID namespace, so
the \texttt{privileged} setting did not help. We did not rule out other causes
experimentally. This is a property of our
Compose configuration, not of Docker in general. We did not test whether
sharing the host PID namespace would allow auditd to run.

Two techniques also left evidence in other log files
(Table~\ref{tab:sources}). Cron logging was off by default in the image and
was enabled by hand at run time (Section~\ref{sec:repro}).

\begin{table}[H]
    \centering
    \caption{Log sources in the lab as configured}
    \label{tab:sources}
    \begin{tabular}{lll}
        \toprule
        \textbf{Source} & \textbf{Status} & \textbf{Techniques with evidence} \\
        \midrule
        \texttt{auth.log} & Available & T1136.001 \\
        \texttt{cron.log} & Available after manual setup & T1053.003 \\
        auditd & Not running & None (9 techniques not evaluated) \\
        \bottomrule
    \end{tabular}
\end{table}

\subsection{Evidence Query}
\label{sec:query}

For each technique the harness ran a keyword query: a time-range filter on
\texttt{@timestamp} combined with a disjunction of \texttt{match} clauses on
the \texttt{message} field. Listing~\ref{lst:query} shows the query for
T1136.001. The keywords for T1053.003 were \texttt{cron\_mod} and
\texttt{crontab}. A technique was counted as observed if the query returned at
least one document. The query was run after the fact; no scheduled rule
produced an alert.

\begin{lstlisting}[caption={Evidence query for T1136.001},label={lst:query}]
{
  "bool": {
    "must": [
      {"range": {"@timestamp": {"gte": "start", "lte": "end"}}},
      {"bool": {
        "should": [
          {"match": {"message": "user_add"}},
          {"match": {"message": "user_mod"}},
          {"match": {"message": "useradd"}}
        ],
        "minimum_should_match": 1
      }}
    ]
  }
}
\end{lstlisting}

\subsection{Collection Windows}

The window for each test started 1\,s before the attack command started and
ended a fixed buffer after it finished. The default buffer was 15\,s. After
the failure described in Section~\ref{sec:cron}, the buffer for T1053.003 was
set to 65\,s, because cron checks the modification time of
\texttt{/etc/crontab} once per minute \cite{cron-man}, so the reload event can
lag the change by up to about 60\,s before shipping delay is added.

\subsection{Measures}

\begin{itemize}
    \item \textbf{Observed}: whether at least one document matched within the
    window.
    \item \textbf{Log-availability latency}: the \texttt{@timestamp} of the
    first matching document at or after the attack start, minus the attack
    start time. Earlier drafts called this ``time to detect.'' It is not a
    detection time, because nothing detected anything. The two timestamps come
    from different clocks: the attack start from the harness on the host, and
    \texttt{@timestamp} from Filebeat, which stamps each line when it reads
    it. Clock offset between the two was not measured.
    \item \textbf{Match count}: the number of matching documents in the window.
\end{itemize}

The harness contains code for a benign baseline and a false-positive rate, but
no baseline was run. Generated reports print a false-positive rate of
0.0 per hour in that case; that value means ``not measured.''

\section{Results}

All runs took place on 2026-10-10. Each result below comes from one run.

\subsection{Pilot Runs of the Full Manifest}

Four runs of the full manifest (12 tests across 11 techniques) between 10:16
and 11:07 local time each reported 0 of 12 tests observed, including the two
techniques that were later observed in single-technique runs. The result files
do not record the cause. The timestamp-handling tests added to the test suite
afterwards suggest an error in converting local time to UTC when building the
query window, but we did not confirm this from the artifacts. Cron logging was
also not yet enabled during these runs. We report the pilot runs because they
are the only runs in which the other nine techniques were executed, and those
nine were not run again after the harness was changed. Their status in
Table~\ref{tab:status} therefore rests on the absence of auditd, not on a
later run that returned no matches.

\subsection{T1136.001: Local Account Creation}

\begin{table}[H]
    \centering
    \caption{T1136.001, single run}
    \label{tab:t1136}
    \begin{tabular}{ll}
        \toprule
        \textbf{Item} & \textbf{Value} \\
        \midrule
        Attack command & \texttt{useradd -M -N -r -s /bin/bash evil\_user} \\
        Command succeeded & Yes \\
        Collection window & 1\,s before start to 15\,s after end \\
        Observed & Yes (1 of 1 run) \\
        Log-availability latency & 3.20\,s \\
        \bottomrule
    \end{tabular}
\end{table}

Evidence file: \texttt{results/validation\_20261010\_125751.json}.

Table~\ref{tab:t1136} gives the run details. The matching event came from
\texttt{auth.log}. The latency of 3.20\,s
covers log writing, Filebeat pickup, and indexing; we did not separate these
components. A separate diagnostic query that added the username
\texttt{evil\_user} returned one document, which supports attributing the
match to this run. That username filter is not part of the harness query, so
the harness alone would also have matched any other \texttt{useradd} activity
in the window. Cleanup (\texttt{userdel}) was run by the harness and was not
checked independently.

\subsection{T1053.003: Cron Persistence}
\label{sec:cron}

\paragraph{Run with the 15\,s window (not observed).}
The attack command appended a line to \texttt{/etc/crontab} at 07:46:37 UTC.
The window closed at 07:46:52 UTC with 0 matching documents
(Table~\ref{tab:cron15}). A wider search afterwards found 2 documents. The one
recorded in the evidence file was indexed at 07:49:44 UTC and read
``\texttt{message repeated 2 times: [ (*system*) RELOAD (/etc/crontab)]}''
with a syslog time of 07:49:01. The reload time of 07:47:01 in the evidence
file is an estimate, not a logged value.

\begin{table}[H]
    \centering
    \caption{T1053.003 with a 15\,s post-attack window, single run}
    \label{tab:cron15}
    \begin{tabular}{ll}
        \toprule
        \textbf{Item} & \textbf{Value} \\
        \midrule
        Attack start & 07:46:37 UTC \\
        Window end & 07:46:52 UTC \\
        Matching documents in window & 0 \\
        Matching documents in wider search & 2 \\
        Observed & No (0 of 1 run) \\
        \bottomrule
    \end{tabular}
\end{table}

Evidence file: \texttt{results/t1053003\_validation\_20261010\_132113.json}.

\paragraph{Run with the 65\,s window (observed).}
With the buffer raised to 65\,s, a second run was observed
(Table~\ref{tab:cron65}).

\begin{table}[H]
    \centering
    \caption{T1053.003 with a 65\,s post-attack window, single run}
    \label{tab:cron65}
    \begin{tabular}{lp{0.62\textwidth}}
        \toprule
        \textbf{Item} & \textbf{Value} \\
        \midrule
        Attack command & \texttt{echo "* * * * * root /tmp/malicious.sh" >{}> /etc/crontab} \\
        Command succeeded & Yes \\
        Attack start & 08:29:33 UTC \\
        Collection window & 1\,s before start to 65\,s after end \\
        Observed & Yes (1 of 1 run) \\
        Matching documents & 2 \\
        Log-availability latency & 30.94\,s \\
        Cleanup checked independently & Yes \\
        \bottomrule
    \end{tabular}
\end{table}

Evidence file: \texttt{results/t1053003\_65s\_window\_20261010\_140043.json}.

Four qualifications apply to this result.

\begin{enumerate}
    \item \textbf{The latency mostly reflects when the command ran.} The
    attack started 33\,s into the minute and cron reloads at the next minute
    boundary, so about 27\,s of the 30.94\,s is waiting for cron. A run
    started at a different second would give a different value. The figure is
    one sample, not a characteristic latency for the technique.
    \item \textbf{The match cannot be tied to the malicious entry.} The
    harness ran the cleanup command immediately after the attack command, so
    the crontab was edited twice before cron next checked it. The matched
    event is a generic reload message that contains the token
    \texttt{crontab}; it does not name the entry that changed.
    \item \textbf{The match count of 2 is one event counted twice.} The
    evidence file stores the count but not the matched documents. The
    documents still held in Elasticsearch, read on 2026-10-13, are the same
    reload line indexed once from \texttt{syslog} and once from
    \texttt{cron.log}, both of which Filebeat ships. The match count is
    therefore inflated for any message written to more than one shipped
    file.
    \item \textbf{A 65\,s window is not a guaranteed bound.} The image sets
    \texttt{\$RepeatedMsgReduction on}, so rsyslog can replace repeated
    identical messages with a later ``message repeated'' line. In the 15\,s
    run this delayed the indexed document to about three minutes after the
    attack.
\end{enumerate}

\subsection{Status of All Planned Techniques}

Table~\ref{tab:status} summarizes the outcome. Evidence was observed for 2 of
11 planned techniques (18\%). Both techniques whose telemetry was available
were observed, in one run each. With one success in one trial, a 95\% Wilson
interval for the per-run probability of observation runs from about 0.21 to 1
\cite{brown2001}, so these runs do not establish that the evidence appears
reliably.

\begin{table}[H]
    \centering
    \caption{Status of the 11 planned techniques. ``Not evaluated'' means the
    technique was executed only in the pilot runs, before the harness was
    changed, and was not run again.}
    \label{tab:status}
    \begin{tabular}{llll}
        \toprule
        \textbf{ID} & \textbf{Technique} & \textbf{Status} & \textbf{Basis} \\
        \midrule
        T1136.001 & Local Account & Observed (1 run) & \texttt{auth.log} \\
        T1053.003 & Cron & Observed (1 run, 65\,s window) & \texttt{cron.log} \\
        \midrule
        T1059.004 & Unix Shell & Not evaluated & auditd not running \\
        T1548.001 & Setuid and Setgid & Not evaluated & auditd not running \\
        T1070.004 & File Deletion & Not evaluated & auditd not running \\
        T1070.003 & Clear Command History & Not evaluated & auditd not running \\
        T1027 & Obfuscated Files or Information & Not evaluated & auditd not running \\
        T1018 & Remote System Discovery & Not evaluated & auditd not running \\
        T1082 & System Information Discovery & Not evaluated & auditd not running \\
        T1105 & Ingress Tool Transfer & Not evaluated & auditd not running \\
        T1571 & Non-Standard Port & Not evaluated & auditd not running \\
        \bottomrule
    \end{tabular}
\end{table}

\section{Threats to Validity and Limitations}
\label{sec:threats}

\paragraph{Construct validity.}
The measurements describe log availability, not detection. Keyword matching on
the message text is weaker than a detection rule: for T1136.001 the keyword
\texttt{useradd} matches the name of the program that wrote the log line. No
rule schedule was involved, so the latencies exclude the delay a scheduled
detection rule would add. Because no benign baseline was run, we cannot say
how often the same keywords match during normal activity; some keywords for
the unobserved techniques (for example \texttt{hostname} and \texttt{nc}) are
common words in log text.

\paragraph{Internal validity.}
The harness and Filebeat clocks were not compared. The cron result is
confounded by the cleanup edit and affected by rsyslog message reduction
(Section~\ref{sec:cron}). The pilot runs failed for a reason that the
artifacts do not record. The two evidence files for the observed techniques
were assembled outside the repository's scripts, so the exact procedure that
produced them is not captured in code. A read-only inspection of the running
target container on 2026-10-13 found two \texttt{rsyslogd} processes. Neither
setup command on record reproduces the container's configuration: the command
in \texttt{results/t1053003\_telemetry\_enabled.json} would leave a malformed
selector in the rsyslog configuration, and the command in earlier drafts leaves
no log path and, in a disposable container, no second daemon. The procedure
actually used to enable cron logging on 2026-10-10 is therefore unknown, and we
do not know whether the second daemon affected the results.
% AUTHOR: confirm you are content to report the 2026-10-13 inspection.

\paragraph{External validity.}
The lab was one Linux container run under Colima on macOS, with one log
store. Nothing here generalizes to Windows hosts, virtual machines, other
shippers, or production log volumes. The nine unobserved techniques may well
be observable in an environment where auditd runs; we did not test one.

\paragraph{Statistical conclusion validity.}
Each technique was run once. No rates, means, or distributions can be
estimated.

\paragraph{Other limitations.}
The Sigma rule files in the repository were not used, and there is no rule
file for T1027. Cron logging required manual setup that does not persist when
the container is recreated. The project directory was not under version
control when the runs were made, so the code state is identified only by file
timestamps.

\section{Discussion}

The answer to the research question is narrow. In this lab, two log sources,
\texttt{auth.log} and \texttt{cron.log}, yielded evidence, for one technique
each. The source that the design depended on, auditd, yielded none; the
likely reason could have been found in the kernel source before the lab was
built.

The cron case shows that a single collection window for all techniques is
inadequate. The delay between an action and its log event depends on the
mechanism that writes the event: \texttt{useradd} logs as it runs, whereas
cron logs a reload at its next minute check. A harness that waits a fixed
15\,s will report a miss for evidence that arrives later, and the miss will
look like a detection failure. Message reduction in the log daemon can extend
the delay further and unpredictably.

Three practices would have prevented most of the problems reported above:
checking that each expected log source produces events before running attacks;
recording the matched documents, not only their count, with each result; and
separating attack and cleanup by more than the longest expected logging delay.
These are suggestions drawn from one small lab, not tested recommendations.

\section{Reproducibility}
\label{sec:repro}

\subsection{Environment}

\begin{itemize}
    \item Host: macOS (Darwin 27.0.0); Python 3.13.5.
    \item Container runtime: Colima, with an Ubuntu 24.04 virtual machine on
    Linux kernel 6.8.0-117, as reported by \texttt{docker info} on 2026-10-13.
    The runtime was not recorded on the day of the runs.
    \item Images: Elasticsearch 8.12.0, Kibana 8.12.0, Filebeat 8.12.0,
    Ubuntu 22.04 (target, built from \texttt{docker/target/Dockerfile}).
\end{itemize}

\subsection{Commands}

Commands are run from the repository root. The first three correspond to
Makefile targets in the repository.

\begin{lstlisting}[language=bash,caption={Install, start the lab, and run the unit tests}]
make install     # pip install -e ".[dev]"
make up          # docker compose -f docker/docker-compose.yml up -d --build
make test        # pytest tests/ -v
\end{lstlisting}

Cron logging must be enabled in the target container before T1053.003 is run.
The command in earlier drafts does not do this correctly:
its \texttt{sed} expression replaces the whole line and removes the log file
path, and \texttt{killall} is not installed in the image. The form below
uncomments the line and uses \texttt{pkill}, which is installed. We
subsequently tested the \texttt{sed} expression and the complete command
sequence in a fresh disposable container created from the
\texttt{docker-target} image; the sequence created
\texttt{/var/log/cron.log}, which contained cron startup lines. The full
\texttt{docker exec soc-target} command was not executed against the live lab
container, so that form remains unverified.

\begin{lstlisting}[language=bash,caption={Enable cron logging in the target container}]
docker exec soc-target bash -c '
  sed -i "s/^#cron\.\*/cron.*/" /etc/rsyslog.d/50-default.conf
  pkill -x rsyslogd; sleep 1; rsyslogd
  service cron start'
\end{lstlisting}

The repository has no script that writes the per-technique evidence files.
The following uses the harness classes that compute the reported measures; the
sleep must exceed the post-attack buffer (15\,s by default, 65\,s for
T1053.003). We verified the class and method signatures but did not execute
this snippet, since doing so would be a new run.

\begin{lstlisting}[language=bash,caption={Run one technique and compute its measures}]
python - <<'EOF'
import time
from benchmark.runner import AttackRunner
from benchmark.collector import AlertCollector
from benchmark.metrics import MetricsCalculator

tid = "T1136.001"
results = AttackRunner().run_single_technique(tid)
time.sleep(20)   # use 70 for T1053.003
print(MetricsCalculator(AlertCollector()).calculate_technique_metrics(tid, results))
EOF
\end{lstlisting}

\subsection{Test Suite}

The suite contains 60 tests, all of which passed when re-run on 2026-10-13.
They replace Elasticsearch with mocks and cover timestamp conversion,
collection-window buffers, keyword matching logic, measure calculation, and
result serialization. They test the harness code; they do not test whether
evidence appears in the lab.

\subsection{Artifacts}

\begin{itemize}
    \item \texttt{results/validation\_20261010\_125751.json}: T1136.001,
    observed.
    \item \texttt{results/t1053003\_65s\_window\_20261010\_140043.json}:
    T1053.003, 65\,s window, observed.
    \item \texttt{results/t1053003\_validation\_20261010\_132113.json}:
    T1053.003, 15\,s window, not observed.
    \item \texttt{results/t1053003\_telemetry\_blocker.json} and
    \texttt{results/t1053003\_telemetry\_enabled.json}: cron logging before
    and after manual setup.
    \item \texttt{results/report\_20261010\_*.json}: the four pilot runs.
\end{itemize}

The repository has not been published at a public address.
% AUTHOR: add the repository URL and a commit or archive identifier when available.

\section{Conclusion}

We set out to benchmark detection rules in a containerized lab and instead
measured whether log evidence arrived at all. Evidence was observed for 2 of
11 planned techniques, in one run each: local account creation after 3.20\,s
and cron persistence after 30.94\,s, the latter only once the post-attack
window was extended from 15\,s to 65\,s. The other 9 techniques were not
evaluated after the harness was corrected, because the audit daemon did not
run in the container; we attribute this to the kernel accepting audit control
only from the initial PID namespace. No detection rule was executed, no
false-positive baseline was
measured, and no run was repeated. The useful results are the reasons for
these outcomes: telemetry preconditions should be verified before attacks are
run, collection windows should follow the logging mechanism of each technique,
and each result should store the matched events so that attribution can be
checked.

\section*{Declarations}

\paragraph{Use of AI tools.}
AI assistance was used for the literature search, for checking the claims in
this report against the repository, and for editing the text.
% AUTHOR: after you have reviewed the paper, add your own statement of review
% and responsibility here. It was removed because it must be made by you.
% AUTHOR: adjust this statement to match your actual use and the venue's policy.
% AUTHOR: add funding, conflict-of-interest, and author-contribution statements
% as the venue requires. They were not drafted because the facts are yours to state.

% ---------------------------------------------------------------------------
% References. Entries are grouped: peer-reviewed research; preprints;
% technical reports, documentation, and source code. Volume and page details
% were not verified and are omitted. AUTHOR: read each source before
% submission; see the revision audit for which summaries rest on abstracts only.
% ---------------------------------------------------------------------------
\begin{thebibliography}{99}

% --- Peer-reviewed research ---

\bibitem{alahmadi2022}
B.~A. Alahmadi, L.~Axon, and I.~Martinovic.
\newblock 99\% False Positives: A Qualitative Study of {SOC} Analysts'
  Perspectives on Security Alarms.
\newblock In \emph{31st USENIX Security Symposium}, 2022.
\newblock \url{https://www.usenix.org/conference/usenixsecurity22/presentation/alahmadi}

\bibitem{applebaum2016}
A.~Applebaum, D.~Miller, B.~Strom, C.~Korban, and R.~Wolf.
\newblock Intelligent, Automated Red Team Emulation.
\newblock In \emph{Annual Computer Security Applications Conference (ACSAC)},
  2016.
\newblock \url{https://doi.org/10.1145/2991079.2991111}

\bibitem{axelsson2000}
S.~Axelsson.
\newblock The Base-Rate Fallacy and the Difficulty of Intrusion Detection.
\newblock \emph{ACM Transactions on Information and System Security}, 2000.
\newblock \url{https://doi.org/10.1145/357830.357849}

\bibitem{brown2001}
L.~D. Brown, T.~T. Cai, and A.~DasGupta.
\newblock Interval Estimation for a Binomial Proportion.
\newblock \emph{Statistical Science}, 2001.
\newblock \url{https://doi.org/10.1214/ss/1009213286}

\bibitem{chen2021}
X.~Chen, H.~Irshad, Y.~Chen, A.~Gehani, and V.~Yegneswaran.
\newblock {CLARION}: Sound and Clear Provenance Tracking for Microservice
  Deployments.
\newblock In \emph{30th USENIX Security Symposium}, 2021.
\newblock \url{https://www.usenix.org/conference/usenixsecurity21/presentation/chen-xutong}

\bibitem{collberg2016}
C.~Collberg and T.~A. Proebsting.
\newblock Repeatability in Computer Systems Research.
\newblock \emph{Communications of the ACM}, 2016.
\newblock \url{https://doi.org/10.1145/2812803}

\bibitem{kokulu2019}
F.~B. Kokulu, A.~Soneji, T.~Bao, Y.~Shoshitaishvili, Z.~Zhao, A.~Doup\'e, and
  G.-J. Ahn.
\newblock Matched and Mismatched {SOCs}: A Qualitative Study on Security
  Operations Center Issues.
\newblock In \emph{ACM SIGSAC Conference on Computer and Communications
  Security (CCS)}, 2019.
\newblock \url{https://doi.org/10.1145/3319535.3354239}

\bibitem{landauer2021}
M.~Landauer, F.~Skopik, M.~Wurzenberger, W.~Hotwagner, and A.~Rauber.
\newblock Have it Your Way: Generating Customized Log Datasets With a
  Model-Driven Simulation Testbed.
\newblock \emph{IEEE Transactions on Reliability}, published online 2020.
\newblock \url{https://doi.org/10.1109/TR.2020.3031317}

\bibitem{landauer2022}
M.~Landauer, F.~Skopik, M.~Frank, W.~Hotwagner, M.~Wurzenberger, and
  A.~Rauber.
\newblock Maintainable Log Datasets for Evaluation of Intrusion Detection
  Systems.
\newblock \emph{IEEE Transactions on Dependable and Secure Computing},
  published online 2022.
\newblock \url{https://doi.org/10.1109/TDSC.2022.3201582}

\bibitem{lim2021}
S.~Y. Lim, B.~Stelea, X.~Han, and T.~Pasquier.
\newblock Secure Namespaced Kernel Audit for Containers.
\newblock In \emph{ACM Symposium on Cloud Computing (SoCC)}, 2021.
\newblock \url{https://doi.org/10.1145/3472883.3486976}

\bibitem{mchugh2000}
J.~McHugh.
\newblock Testing Intrusion Detection Systems: A Critique of the 1998 and 1999
  {DARPA} Intrusion Detection System Evaluations as Performed by Lincoln
  Laboratory.
\newblock \emph{ACM Transactions on Information and System Security}, 2000.
\newblock \url{https://doi.org/10.1145/382912.382923}

\bibitem{shen2024}
X.~Shen, Z.~Li, G.~Burleigh, L.~Wang, and Y.~Chen.
\newblock Decoding the {MITRE} Engenuity {ATT\&CK} Enterprise Evaluation: An
  Analysis of {EDR} Performance in Real-World Environments.
\newblock In \emph{ACM Asia Conference on Computer and Communications Security
  (ASIA CCS)}, 2024.
\newblock \url{https://doi.org/10.1145/3634737.3645012}

\bibitem{uetz2021}
R.~Uetz, C.~Hemminghaus, L.~Hackl\"ander, P.~Schlipper, and M.~Henze.
\newblock Reproducible and Adaptable Log Data Generation for Sound
  Cybersecurity Experiments.
\newblock In \emph{Annual Computer Security Applications Conference (ACSAC)},
  2021.
\newblock \url{https://doi.org/10.1145/3485832.3488020}

\bibitem{uetz2024}
R.~Uetz, M.~Herzog, L.~Hackl\"ander, S.~Schwarz, and M.~Henze.
\newblock You Cannot Escape Me: Detecting Evasions of {SIEM} Rules in
  Enterprise Networks.
\newblock In \emph{33rd USENIX Security Symposium}, 2024.
\newblock \url{https://www.usenix.org/conference/usenixsecurity24/presentation/uetz}

\bibitem{vanderkouwe2019}
E.~van~der Kouwe, G.~Heiser, D.~Andriesse, H.~Bos, and C.~Giuffrida.
\newblock {SoK}: Benchmarking Flaws in Systems Security.
\newblock In \emph{IEEE European Symposium on Security and Privacy
  (EuroS\&P)}, 2019.
\newblock \url{https://doi.org/10.1109/EuroSP.2019.00031}

\bibitem{vermeer2022}
M.~Vermeer, M.~van Eeten, and C.~Ga\~n\'an.
\newblock Ruling the Rules: Quantifying the Evolution of Rulesets, Alerts and
  Incidents in Network Intrusion Detection.
\newblock In \emph{ACM Asia Conference on Computer and Communications Security
  (ASIA CCS)}, 2022.
\newblock \url{https://doi.org/10.1145/3488932.3517412}

\bibitem{vermeer2023}
M.~Vermeer, N.~Kadenko, M.~van Eeten, C.~Ga\~n\'an, and S.~Parkin.
\newblock Alert Alchemy: {SOC} Workflows and Decisions in the Management of
  {NIDS} Rules.
\newblock In \emph{ACM SIGSAC Conference on Computer and Communications
  Security (CCS)}, 2023.
\newblock \url{https://doi.org/10.1145/3576915.3616581}

\bibitem{virkud2024}
A.~Virkud, M.~A. Inam, A.~Riddle, J.~Liu, G.~Wang, and A.~Bates.
\newblock How does Endpoint Detection use the {MITRE ATT\&CK} Framework?
\newblock In \emph{33rd USENIX Security Symposium}, 2024.
\newblock \url{https://www.usenix.org/conference/usenixsecurity24/presentation/virkud}

% --- Preprints ---

\bibitem{houssel2026}
P.~R.~B. Houssel, O.~Levillain, S.~Laniepce, N.~Dejon, and H.~Debar.
\newblock A Study of Kernel Telemetry Options for Security-Oriented
  Provenance.
\newblock arXiv:2608.11418, 2026. Preprint; the arXiv record states it was
  accepted at ISC 2026.
\newblock \url{https://arxiv.org/abs/2608.11418}

\bibitem{long2026}
M.~Long and D.~Evans.
\newblock Evolution of Log-Based Detection Rules in Public Repositories.
\newblock arXiv:2605.05383, 2026. Preprint, not peer reviewed.
\newblock \url{https://arxiv.org/abs/2605.05383}

\bibitem{shukla2025}
A.~Shukla, P.~A. Gandhi, Y.~Elovici, and A.~Shabtai.
\newblock {RuleGenie}: {SIEM} Detection Rule Set Optimization.
\newblock arXiv:2505.06701, 2025. Preprint, not peer reviewed.
\newblock \url{https://arxiv.org/abs/2505.06701}

% --- Technical reports, documentation, and source code ---

\bibitem{attack-design}
The MITRE Corporation.
\newblock {MITRE ATT\&CK}: Design and Philosophy.
\newblock Technical report, revised March 2020.
\newblock \url{https://attack.mitre.org/docs/ATTACK_Design_and_Philosophy_March_2020.pdf}

\bibitem{attack-evals}
The MITRE Corporation.
\newblock {ATT\&CK} Evaluations.
\newblock Project website.
\newblock \url{https://evals.mitre.org/}

\bibitem{atomic}
Red Canary.
\newblock Atomic Red Team.
\newblock Software repository and documentation.
\newblock \url{https://github.com/redcanaryco/atomic-red-team}

\bibitem{cron-man}
Ubuntu Manpage Repository.
\newblock cron(8), Ubuntu 22.04 (jammy).
\newblock Documentation.
\newblock \url{https://manpages.ubuntu.com/manpages/jammy/man8/cron.8.html}

\bibitem{linux-audit}
Linux kernel developers.
\newblock \texttt{kernel/audit.c}, function \texttt{audit\_netlink\_ok}.
\newblock Source code.
\newblock \url{https://github.com/torvalds/linux/blob/master/kernel/audit.c}

\bibitem{pysigma}
SigmaHQ.
\newblock pySigma.
\newblock Software repository.
\newblock \url{https://github.com/SigmaHQ/pySigma}

\bibitem{sigma-spec}
SigmaHQ.
\newblock Sigma Specification: Generic Signature Format for {SIEM} Systems.
\newblock Specification repository.
\newblock \url{https://github.com/SigmaHQ/sigma-specification}

\end{thebibliography}

\end{document}
